\documentclass[a4j]{jarticle}
\usepackage{ascmac}
\usepackage[dvipdfmx]{graphicx}
\usepackage{listings,jvlisting}

\lstset{
  basicstyle={\ttfamily},
  identifierstyle={\small},
  commentstyle={\smallitshape},
  keywordstyle={\small\bfseries},
  ndkeywordstyle={\small},
  stringstyle={\small\ttfamily},
  frame={tb},
  breaklines=true,
  columns=[l]{fullflexible},
  numbers=left,
  xrightmargin=0zw,
  xleftmargin=3zw,
  numberstyle={\scriptsize},
  stepnumber=1,
  numbersep=1zw,
  lineskip=-0.5ex
}

\begin{document}

\vfill
\begin{screen}
  \begin{center}
    \huge\bf 情報科学概論　第十回レポート課題\\
    \vspace{5mm}
    \large\bf
    学籍番号　2022311042 \quad 名前 神野友哉\\
    \today
  \end{center}
\end{screen}
\thispagestyle{empty}

\section{Python 入門実習　問１}
第七回授業のMatlabコード(contr4pend.m)をpythonコードに変換せよ。同等の
シミュレーションであることを示す結果と比較検討を示すこと。(アニメーション部分は任意)
\subsection{結果}
図 \ref{fig:e}はコード \ref{python}(python)の実行結果であり、図 \ref{fig:f}も同様にコード \ref{matlab}(matlab)の
実行結果である。図 \ref{fig:e}と図 \ref{fig:f}は同じ結果となった。言語仕様に困った点としては、
pythonの常微分方程式を解く関数odeintと、matlabの同様の関数ode45は出力において、それぞれxのみか、xとその対応するtが纏まって
行列として出るかの違いがあった。さらに、行列のインデックスの始まりが異なり、pythonでは0から始まるが、matlabでは1から始まっている。
行列の積においても、pythonはnp.dot(a, b)のように書く必要があるが、matlabでは*演算子のみで書ける。
matlabの @(仮引数) 関数 で書かれるラムダ式の、pythonでの書き方が余り理解できなかった。

\begin{figure}[htbp]
  \begin{minipage}{0.5\hsize}
    \begin{center}
      \includegraphics[width=70mm]{python-result.png}
    \end{center}
    \caption{pythonの結果}
    \label{fig:e}
  \end{minipage}
  \begin{minipage}{0.5\hsize}
    \begin{center}
      \includegraphics[width=70mm]{matlab-result.png}
    \end{center}
    \caption{matlabの結果}
    \label{fig:f}
  \end{minipage}
\end{figure}

\begin{lstlisting}[caption=python,label=python]
import numpy as np
from scipy.integrate import odeint
import matplotlib.pyplot as plt

# 物理パラメータ
M = 1         # 質量 [kg]
l = 0.5       # 回転軸から重心までの距離 [m]
J = 1         # 慣性モーメント [kg m^2]
c = 0.42      # 粘性摩擦係数 [Nm/(rad/s)] ※XX : 学籍番号の下 2 桁
ga = 0     # 重力加速度 [m/s^2]

# 時間パラメータ
t0 = 0
dt = 0.01
tf = 10

# 初期状態 9 初期時刻での状態)
x0 = [np.pi/2, 0]    # x(0) = [ theta(0), thetadot(0) ]

# 制御ゲイン
kp = 1   # 比例(Proportional; P)ゲイン
kd = 1   # 微分(Differential; D)ゲイン

# 各種制御則
def ux(t, x):
    # 1) 無制御（駆動なし）
    #u = 0*t
    # 2) PD 制御（ただし, kd=0 => P 制御, kp=0 => D 制御）
    u = -np.dot([kp, kd], x)
    return u

# 詳細モデル（非線形システム）
def pend_nlnsys(x, t):
    theta, thetadot = x
    # システム表現の一部（ベクトル場）
    f = np.array([thetadot, -M*ga*l/J*np.sin(theta)-c/J*thetadot])
    g = np.array([0, 1/J])
    # ODE
    return f + np.dot(g, ux(t, x))

# 設計用モデル（x = xe まわりで近似した線形システム）
A = np.array([[0, 1], [-M*ga*l/J, -c/J]])
b = np.array([0, 1/J])
print("A の固有値は ")
print(np.linalg.eigvals(A))

def pend_alnsys(x, t):
    return np.dot(A, x) + b * ux(t, x)

# 状態の解軌道を数値計算
t4nls = np.arange(t0, tf+dt, dt)
x4nls = odeint(pend_nlnsys, x0, t4nls)

# 制御入力の時系列データ算出
u4nls = [ux(t, x) for t, x in zip(t4nls, x4nls)]


# 状態の解軌道を数値計算（設計用モデル）
t4als = np.arange(t0, tf+dt, dt)
x4als = odeint(pend_alnsys, x0, t4als)

# 制御入力の時系列データ算出
u4als = [ux(t, x) for t, x in zip(t4als, x4als)]

# 結果表示
plt.figure(1)

plt.subplot(3, 1, 1)
plt.plot(t4nls, x4nls[:, 0], 'b-', linewidth=2, label='original')
plt.plot(t4als, x4als[:, 0], 'r--', linewidth=2, label='linearized')
plt.legend(loc='best')
plt.ylabel(r'$\theta$ [rad]')
plt.grid(True)

plt.subplot(3, 1, 2)
plt.plot(t4nls, x4nls[:, 1], 'b-', linewidth=2, label='original')
plt.plot(t4als, x4als[:, 1], 'r--', linewidth=2, label='linearized')
plt.legend(loc='best')
plt.ylabel(r'$\frac{d\theta}{dt}$ [rad/s]')
plt.grid(True)

plt.subplot(3, 1, 3)
plt.plot(t4nls, u4nls, 'b-', linewidth=2, label='for original')
plt.plot(t4als, u4als, 'r--', linewidth=2, label='for linearized')
plt.legend(loc='best')
plt.xlabel('t [s]')
plt.ylabel(r'$\tau$ [Nm]')
plt.grid(True)

plt.show()
\end{lstlisting}

\begin{lstlisting}[caption=matlab,label=matlab]
close all         % すべてのFigure を閉じる
clear variables   % すべての変数をクリアする
  
%% パラメータ設定
% 物理パラメータ
global M l J c ga
M  = 1;      % 質量 [kg]
l  = 0.5;    % 回転軸から重心までの距離 [m]
J  = 1;      % 慣性モーメント [kg m^2]
c  = 0.42;   % 粘性摩擦係数 [Nm/(rad/s)] ※XX : 学籍番号の下 2 桁
ga = 0;   % 重力加速度 [m/s^2]
  
  
% 時間パラメータ
t0 = 0;  dt = 0.01;  tf = 10;
  
% 初期状態（初期時刻での状態）
x0 = [ -pi/2, 0 ];    %x(0) = [ theta(0), thetadot(0) ]
  
% 制御ゲイン（パラメータ
global kp kd
kp = 1;   % 比例（Proportional; P）ゲイン
kd = 1;   % 微分（Differential; D）ゲイン
  
  
  
%% 常微分方程式を数値的に解く：詳細モデル（非線形システム）の場合
% 状態の解軌道を数値計算（nls: non-linear system の略）
[t4nls, x4nls] = ode45( @(t,x) pend_nlnsys(t,x), t0:dt:tf, x0 );
% 制御入力の時系列データ算出
u4nls = arrayfun(@(t,x1,x2) ux(t,[x1;x2]), t4nls,x4nls(:,1),x4nls(:,2));
  
%% 常微分方程式を数値的に解く: 設計用モデル（x = xe まわりで近似した線形システム）の場合
% システム表現の一部（システム行列
global A b
A = [0,1;-M*ga*l/J,-c/J];   % 適切に定義する
b = [0;1/J];   % 適切に定義する
disp('A の固有値は')
eig(A)
% 状態の解軌道を数値計算（als: approximately linearized system の略）
[t4als, x4als] = ode45( @(t,x) pend_alnsys(t,x), t0:dt:tf, x0 );
% 制御入力の時系列データ算出
u4als = arrayfun(@(t,x1,x2) ux(t,[x1;x2]), t4als,x4als(:,1),x4als(:,2));
  
%% 結果表示（データの可視化）
figure(1)
subplot(3,1,1)
plot( t4nls,x4nls(:,1), 'LineWidth',2), hold on
plot( t4als,x4als(:,1), '--','LineWidth',2)
legend('original','linearized', 'Location','best')
ylabel('\theta [rad]'), grid on
subplot(3,1,2)
plot( t4nls,x4nls(:,2), 'LineWidth',2 ), hold on
plot( t4als,x4als(:,2), '--','LineWidth',2)
legend('original','linearized', 'Location','best')
ylabel('d\theta/dt [rad/s]'), grid on
subplot(3,1,3)
plot( t4nls,u4nls, 'LineWidth',2), hold on
plot( t4als,u4als, '--','LineWidth',2)
legend('for original','for linearized', 'Location','best')
xlabel('t [s]'), ylabel('\tau [Nm]'), grid on
print -dpdf -fillpage 'result_t-x+u.pdf'
%print -depsc 'result_t-x+u.eps'
  
show_anim(2,l,t4nls,x4nls)
  
  
%% ローカル関数
  
% 各種制御則（いずれか一行のみ有効化し，必要に応じて編集する）
function u = ux(t, x)
    global kp kd
    % 1) 無制御（駆動なし）
    %u = 0*t;
    % 2) PD 制御（ただし, kd=0 => P 制御, kp=0 => D 制御）
    u = -[kp kd]*x;    %(= -kp*x(1) - kd*x(2))
end
  
% 詳細モデル（非線形システム）
function dxdt = pend_nlnsys(t, x)
    global M l J c ga
    % システム表現の一部（ベクトル場）
    f = [x(2);-M*ga*l/J*sin(x(1))-c/J*x(2)];   % 適切に定義する
    g = [   0;
          1/J ];
    % ODE
    dxdt = f + g*ux(t,x);
end
  
% 設計用モデル（x = xe まわりで近似した線形システム）
function dxdt = pend_alnsys(t, x)
    global A b
    % ODE
    dxdt = A*x + b*ux(t,x);
end
  
function show_anim(fid,l,tc,xc)
    figure(fid), clf
    ax = gca;  ax.XAxisLocation = 'origin';  ax.YAxisLocation = 'origin';
    axis equal, xlim([-2.5*l 2.5*l]), ylim([-2.5*l 2.5*l])
    xlabel('X (m)'), ylabel('Y (m)'), grid on, hold off, hold on
    p1 = @(t) zeros(numel(t), 2);
    pE = @(theta) [ 2*l*sin(theta), -2*l*cos(theta) ];
    alpha = 0.5;  p1t = p1(tc);  pEt = pE(xc(:,1));
    for i=1:(numel(tc)-1)/50:numel(tc)
        if i==1
            plot([p1t(i,1) pEt(i,1)],...
                 [p1t(i,2) pEt(i,2)],'b-', 'Linewidth',5)
            plot(p1t(i,1),p1t(i,2),'bo', 'MarkerSize',10,...
                 'MarkerFaceColor',[0,0,0]+alpha,'MarkerEdgeColor','b')
            plot(pEt(i,1),pEt(i,2),'bo', 'MarkerSize',10,...
                 'MarkerFaceColor','m', 'MarkerEdgeColor','b')
        elseif i==numel(tc)
            plot([p1t(i,1) pEt(i,1)],...
                 [p1t(i,2) pEt(i,2)],'k-', 'Linewidth',5)
            plot(p1t(i,1),p1t(i,2),'ko', 'MarkerSize',10,...
                 'MarkerFaceColor',[0,0,0]+alpha,'MarkerEdgeColor','k')
            plot(pEt(i,1),pEt(i,2),'ko', 'MarkerSize',10,...
                 'MarkerFaceColor','m', 'MarkerEdgeColor','k')
        else
            plot([p1t(i,1) pEt(i,1)],...
                 [p1t(i,2) pEt(i,2)],'c-', 'Linewidth',3)
        end
        pause(.25)
    end
    plot(pEt(:,1),pEt(:,2),'r--', 'Linewidth',3), hold off
    print -dpdf -fillpage 'result_link-motion.pdf'
    %print -depsc 'result_link-motion.eps'
end
\end{lstlisting}
\end{document}